\BourbakiStarSection{Bibliography}

\begin{enumerate}
\def\labelenumi{\arabic{enumi}.}
\item
  O. NEUGEBAUER, Vorlesungen über Geschichte der antiken Mathematik, Vol. I: Vorgriechische Mathematik, Berlin (Springer), 1934.
\item
  Euclidis Elementa, 5 vols., ed.~J. L. Heiberg, Lipsiae (Teubner), 1883--88.
\end{enumerate}

2 bis. T. L. HEATH, The thirteen books of Euclid's Elements . . ., 3 vols., Cambridge, 1908.

\begin{enumerate}
\def\labelenumi{\arabic{enumi}.}
\setcounter{enumi}{2}
\tightlist
\item
  Diophanti Alexandrini Opera Omnia \ldots, 2 vols., ed.~P. Tannery, Lipsiae (Teubner), 1893--95.
\end{enumerate}

3 bis. Diophante d'Alexandrie, trad. P. Ver Eecke, Bruges (Desclée-de Brouwer), 1926.

\begin{enumerate}
\def\labelenumi{\arabic{enumi}.}
\setcounter{enumi}{3}
\item
  G. W. LEIBNIZ, Mathematische Schriften, ed.~C. I. Gerhardt, vol.~V, Halle (Schmidt), 1858.
\item
  C. F. GAUSS, Werke, vols. I (Göttingen, 1870), II (ibid., 1863) and VIII (ibid., 1900).
\item
  A. L. CAUCHY, Œuvres complètes (2), vol.~I, Paris (Gauthier-Villars), 1905.
\item
  N. H. ABEL, Œuvres, 2 vol., ed.~Sylow and Lie, Christiania, 1881.
\item
  E. GALOIS, Écrits et mémoires mathématiques, ed.~R. Bourgne and J. Y. Azra, Paris (Gauthier-Villars), 1962.
\item
  W. R. HAMILTON, Lectures on Quaternions, Dublin, 1853.
\end{enumerate}

9 bis. W. R. HAMILTON, Memorandum respecting a new system of roots of unity, Phil. Mag. (4), 12 (1856), p.~446.

\begin{enumerate}
\def\labelenumi{\arabic{enumi}.}
\setcounter{enumi}{9}
\item
  A. CAYLEY, Collected mathematical papers, vols. I and II, Cambridge (University Press), 1889.
\item
  H. HANKEL, Vorlesungen über die complexen Zahlen und ihre Functionen, 1st part: Theorie der complexen Zahlensysteme, Leipzig (Voss), 1867.
\item
  J. A. SERRET, Cours d'Algèbre supérieure, 3rd ed., Paris (Gauthier-Villars), 1866.
\item
  R. DEDEKIND and H. WEBER, Theorie der algebraischen Funktionen einer Veränderlichen, Crelle's J., 92 (1882), pp.~181--290.
\item
  R. DEDEKIND, Gesammelte mathematische Werke, 3 vols., Braunschweig (Vieweg), 1932.
\item
  C. JORDAN, Traité des substitutions et des équations algébriques, Paris (Gauthier-Villars), 1870 (reprinted, Paris (A. Blanchard), 1957).
\item
  C. JORDAN, Mémoire sur les groupes des mouvements, Ann. di Mat. (2), 2 (1868), pp.~167--215 and 322--345 (= Œuvres, vol.~IV, pp.~231--302, Paris (Gauthier-Villars), 1964).
\item
  L. SYLOW, Théorèmes sur les groupes de substitutions, Math. Ann., 5 (1872), pp.~584--594.
\item
  W. DYCK, Gruppentheoretische Studien, Math. Ann., 20 (1882), pp.~1-44.
\item
  G. FROBENIUS, Über Gruppencharaktere, Berliner Sitzungsber., 1896, pp.~985--1021 (= Gesammelte Abhandlungen, ed.~J. P. Serre, Berlin-Heidelberg-New York (Springer), vol.~III (1968), pp.~1--37).
\item
  W. BURNSIDE, Theory of groups of finite order, 2nd ed., Cambridge, 1911.
\item
  J. NIELSEN, Die Isomorphismengruppen der freien Gruppen, Math. Ann., 91 (1924), pp.~169--209.
\item
  O. SCHREIER, Die Untergruppen der freien Gruppe, Abh. Hamb., 5 (1927), pp.~161--185.
\end{enumerate}

\begin{enumerate}
\def\labelenumi{\arabic{enumi}.}
\setcounter{enumi}{22}
\item
  A. SPEISER, Theorie der Gruppen von endlicher Ordnung, 3rd ed., Berlin (Springer), 1937.
\item
  H. ZASSENHAUS, Lehrbuch der Gruppentheorie, vol.~I, Leipzig-Berlin (Teubner), 1937.
\item
  B. L. VAN DER WAERDEN, Moderne Algebra, 2nd ed., vol.~I, Berlin (Springer), 1937; vol.~II (ibid.), 1940.
\item
  W. MAGNUS, A. KARRASS and D. SOLITAR, Combinatorial group theory, New York (Interscience), 1966.
\item
  D. GORENSTEIN, Finite groups, New York (Harper and Row), 1968.
\end{enumerate}

